\documentclass[11pt]{article}
\usepackage[a4paper,margin=27mm]{geometry}
\usepackage{amsmath,amssymb,amsthm,hyperref}
\newcommand{\PSH}{\operatorname{PSH}}
\newcommand{\FS}{\operatorname{FS}}
\newcommand{\vol}{\operatorname{vol}}
\newtheorem{theorem}{Theorem}
\newtheorem{lemma}[theorem]{Lemma}
\title{Entire big exceptional cones in every dimension}
\author{}
\date{6 October 2026}
\begin{document}
\maketitle
\begin{abstract}
This unreviewed candidate proves a structural theorem for the entire
genuine cone of $L=\pi^*A+cE$, where $\pi$ is the blowup of a
smooth center of codimension at least two, $A$ is actual ample and
$c>0$. The conclusion is in every dimension, on the original
analytic carrier and over every algebraically closed trivially
valued field of characteristic zero. Exact section division gives
the entire cone as the fixed exceptional potential plus pullbacks
of the ample cone. Compactness and arbitrary-net decreasing closure
follow without a general nonample flattening theorem. The theorem
is a narrower structural gate; it does not solve original P16.5's
arbitrary big classes or construct its entire finite-energy equation.
\end{abstract}

\section{The exact domain and objects}
Let $k$ be any algebraically closed trivially valued field of
characteristic zero, $Y$ a smooth integral projective variety of
dimension $n\geq2$, and $Z\subset Y$ a nonempty smooth closed
center of codimension at least two. Let
$\pi:X=\operatorname{Bl}_Z Y\to Y$ and let $E$ be its
effective exceptional Cartier divisor. The source $X$ is smooth
integral projective. Fix an actual ample rational bundle $A$ on $Y$,
a rational $c>0$, and the actual big rational bundle
\[
 P=\pi^*A,\qquad L=P+cE.                                  \tag{1}
\]
The class is nonnef: a line in a geometric exceptional projective
space has $L$-degree $-c$. The center may have several disjoint
components; the single Cartier divisor $E$ and the same coefficient
$c$ are retained. No cardinality hypothesis on $k$ is imposed.

For an actual rational bundle $B$ a generically finite FS potential is
\[
 f=m^{-1}\max_j(\log|s_j|+\lambda_j),\qquad
 0\ne s_j\in H^0(mB),\quad \lambda_j\in\mathbf R.          \tag{2}
\]
Canonical trivial norms are used. The section system need not be
generated, and its poles are allowed. The genuine numerical cone is
exactly the class of one-net limits
\[
 \PSH(\theta)=\{u=\inf_i f_i:\ f_i\searrow\text{ on all points},
  \ f_i\in\FS^{\rm gf}(B_i),\ c_1(B_i)\to\theta,
                         u(x_0)>-\infty\}.               \tag{3}
\]
Here $x_0$ is the generic trivial valuation, and the index set may
have any cardinality. Normalization means $u(x_0)=0$. The topology
is the pointwise topology at divisorial valuations, not the all-point
topology. Birational proper $X,Y$ have the same set of these
valuations of the common function field.

The sole analytic premise $G_0$ is the exact whole actual-ample
theorem retained for original P16.1, applied separately to these
smooth projective $X$ and $Y$: finite continuous same-bundle FS
regularization of genuine ample potentials, decreasing closure with
finite generic endpoint, normalized divisorial compactness, finite
divisorial coordinates and all-point determination. No historical
general-big theorem or disputed abstract extended-real flattening
lemma is a premise. Necessary nef and intersection consequences are
proved below. Ordinary blowup, ample and normal-section facts are
also stated with their applications.

\section{Elementary genuine constructions}
An actual ample rational bundle has a generated power whose
weight-zero FS potential is identically zero. At every valuation-ring
center one generator is a unit and all generators have norm at most
one. Tensoring this system places a function (2) unchanged on the
sum of its bundle and that ample bundle. Hence every already
genuine function in a numerical class $\theta$ embeds, on a tail
of its defining net, in any fixed actual ample bundle $B$ with
$c_1(B)-\theta$ ample. Indeed $B-B_i$ is then actual ample.
This proves all-point determination in every genuine numerical class
by putting two functions in a common ample cone; it does not assert
that a new nonample limit is genuine.

Products of all section monomials after denominator clearing give
literal sums of FS potentials. A maximum of two FS potentials on
one actual bundle is FS there: take all required powers of each
system and then their weighted union. The same maximum is genuine
for two ample potentials, by taking maxima of their two finite
same-bundle regularizing nets, on the product-directed index set.
Consequently divisorial order for two already genuine potentials
in one numerical class implies all-point order: their maximum
in a common ample cone has the larger one's divisorial coordinates,
and $G_0$ gives equality everywhere.

We record the intersection construction needed to handle a limiting
class, with its pole issue resolved explicitly.
\begin{lemma}[Nef intersection]
For actual nef $Q$ and actual ample $H$ on either carrier,
\[
 \PSH(c_1(Q))=\bigcap_{\epsilon\in\mathbf Q_{>0}}
                         \PSH(Q+\epsilon H).             \tag{4}
\]
It is closed under arbitrary decreasing nets with finite limiting
generic value. Its normalized cone is compact in divisorial
coordinates and has one-net finite FS regularization in classes
tending to $c_1(Q)$.
\end{lemma}
\begin{proof}
The forward direction is the zero-system section embedding.
For the reverse direction let $u$ be in every displayed ample cone
and choose finite continuous $f_{\epsilon,j}\searrow u$ on
the actual bundle $Q+\epsilon H$. Consider all triples
$(\epsilon,j,a)$ with $a>0$, ordered by decreasing $\epsilon$,
decreasing $a$ and decreasing function $f_{\epsilon,j}+a$.
To dominate two triples choose smaller $\epsilon_3$ and smaller
positive $a_3$ than both shifts. The finite continuous function
$R=\min(f_{\epsilon_1,j_1}+a_1,f_{\epsilon_2,j_2}+a_2)$
strictly exceeds $u+a_3$. This is also strict at every pole of $u$,
since $R$ is finite. The open sets
\[
 \{\exp(f_{\epsilon_3,j}+a_3)-\exp R<0\}
\]
cover the compact analytic space. A finite subcover and directedness
give a common $j$ for which the comparison holds everywhere.
The triple system is therefore one decreasing FS net. Its classes
tend to $c_1(Q)$; at each point finite thresholds and arbitrarily
small shifts identify its infimum with $u$, also when $u=-\infty$.
This proves (4) in the exact definition (3).

An outer decreasing net belongs unchanged to every ample cone in
(4). Apply $G_0$ there and then (4). For compactness embed every
normalized ample cone $Q+\epsilon H$, $0<\epsilon\leq1$, into
the normalized cone $Q+H$; its compact image is closed in the
Hausdorff coordinate space. Their intersection is compact. All-point
determination in that common ample cone identifies the intersection
with the normalized original representatives in (4).
\end{proof}

For an actual ample $A$ on $Y$ the special intersection
$\bigcap_{\delta>0}\PSH((1+d\delta)A)=\PSH(A)$
is thus valid, using $Q=A$ and $H=dA$. Each member is understood
as the same actual all-point function, with the same normalization.

\section{Actual section division, including perturbations}
The blowup tautological bundle $\mathcal O_X(-E)$ is relatively
ample. Therefore, for a sufficiently large divisible integer $d$,
\[
 H=d\pi^*A-E\quad\hbox{is actual ample on }X.              \tag{5}
\]
Here this is the usual projective relative-ample plus pullback-ample
fact: embed a relatively very ample power into a projective bundle
over $Y$, tensor with a sufficiently large ample power on $Y$,
and use the absolute very ample embedding. Rational denominators
are cleared first.

For every integer $q\geq0$,
\[
                    \pi_*\mathcal O_X(qE)=\mathcal O_Y.  \tag{6}
\]
Indeed a section on $\pi^{-1}U$ is a rational function regular
off $Z\cap U$. It has no pole at any codimension-one point of
the normal variety $U$, since the removed center has codimension
at least two. A normal noetherian domain is the intersection of
its height-one valuation rings inside its fraction field, so it
extends to a regular function on $U$. Conversely its pullback is
regular and hence is a section of $\mathcal O_X(qE)$.
The identification as a section includes multiplication by the
canonical section $s_{qE}$ in local line-bundle frames. This proves
(6) as a sheaf assertion on each affine $U$, and they agree on
overlaps.

The projection formula, or the same local-frame proof, gives for
every actual rational bundle $B$ on $Y$ and rational $b\geq0$
the exact isomorphism
\[
 H^0(Y,mB)\longrightarrow H^0(X,m\pi^*B+mbE),\qquad
                  t\longmapsto s_{mbE}\pi^*t,             \tag{7}
\]
for all sufficiently divisible $m$. Every section on the right is
divisible by the prescribed exceptional factor; no numerical
equivalence is substituted for this actual identity.

For rational $0<\epsilon<c$ put
\[
 A_\epsilon=(1+d\epsilon)A,\quad
 N_\epsilon=(c-\epsilon)E,\quad
 L_\epsilon=L+\epsilon H=\pi^*A_\epsilon+N_\epsilon.
                                                               \tag{8}
\]
These are actual bundles and divisors, $A_\epsilon$ is ample,
$N_\epsilon$ is effective, and (7) applies to every actual section
of every sufficiently divisible power of $L_\epsilon$.
There is no nef or compactness assumption about $L_\epsilon$.
Let $\sigma_b=b\log|s_E|$; denominator clearing makes this an
actual one-section FS potential of $bE$ for rational $b\geq0$.
Its values are finite on divisorial points and zero at $x_0$.

\section{The entire cone theorem}
\begin{theorem}[Unreviewed structural candidate]
In the domain (1), the full genuine numerical cone is
\[
 \PSH(c_1(L))
       =\{\sigma_c+\pi^*w:w\in\PSH(Y,A)\}.              \tag{9}
\]
Every equality in (9) holds at all analytic points of the original
$X^{\rm an}$. On normalized cones this map is a homeomorphism
in divisorial coordinates. In particular the entire genuine big
cone is normalized compact and closed under arbitrary decreasing
nets with finite limiting generic value. Moreover
\[
 \PSH(c_1(\pi^*A))=\pi^*\PSH(Y,A),\qquad
                     \vol(L)=A^n.                         \tag{10}
\]
\end{theorem}
\begin{proof}
If $w\in\PSH(Y,A)$, its finite same-bundle regularizers pull
back to actual FS on $\pi^*A$, with the same decreasing
all-point limit. Tensoring the fixed $s_{mcE}$ gives one defining
FS net for $\sigma_c+\pi^*w$ on $L$ itself. Thus the right
side of (9) is genuine, with exactly the fixed exceptional ideal
powers rather than merely the same pole set.

Conversely let $u\in\PSH(c_1(L))$, normalized. Fix
$0<\epsilon<c$. If $f_i\in\FS^{\rm gf}(B_i)$ is its original
defining net, then $L_\epsilon-B_i$ is actual ample on a tail,
because its numerical class tends to $\epsilon c_1(H)$.
Tensor the zero FS systems to put the unchanged $f_i$ on
$L_\epsilon$. Clear the denominator of $N_\epsilon$ by
all products of each section system; this keeps the function
unchanged. Equation (7) divides every resulting section exactly
and gives actual FS functions $g_i$ on $Y,A_\epsilon$ with
\[
                 f_i=\sigma_{c-\epsilon}+\pi^*g_i        \tag{11}
\]
at all points. At poles this is simply the norm identity for a
product of sections; it uses no extended-real subtraction.

The common function-field divisorial evaluations of (11) are
finite, so $g_i$ decrease at every divisor. Their actual maximum
and $G_0$ on $Y,A_\epsilon$ upgrade this to all-point decrease.
The finite limiting generic value is zero. Thus $G_0$, or directly
the one FS net in (3), gives $w_\epsilon\in\PSH(Y,A_\epsilon)$
as their all-point infimum, normalized. Taking the literal infimum
in (11) yields
\[
 u=\sigma_{c-\epsilon}+\pi^*w_\epsilon
 \quad\hbox{on all of }X^{\rm an}.                         \tag{12}
\]
At points where the exceptional factor is infinite the equality
follows since every $f_i$ is $-\infty$; elsewhere its fixed finite
value commutes with the infimum. At a common divisor $v$ it reads
\[
 w_\epsilon(v)=u(v)-(c-\epsilon)\log|s_E|(v).              \tag{13}
\]

As rational $\epsilon\downarrow0$, the normalized $w_\epsilon$
all embed in a single normalized ample cone on $Y$. Compactness
gives a convergent subnet. For each rational $\delta>0$ the
subnet tail embeds in the compact closed normalized cone of
$(1+d\delta)A$, because $\epsilon<\delta$ eventually.
Its limit belongs to every such cone. The explicitly proved
intersection construction gives an actual $w\in\PSH(Y,A)$.
Equation (13) prescribes every finite divisorial limit as
$u(v)-c\log|s_E|(v)$. The already genuine function
$\sigma_c+\pi^*w$ therefore agrees with $u$ divisorially.
Both embed in one actual ample cone on $X$, where determination
gives equality at all points, including the poles. This proves (9).

On the common divisorial coordinates the map is fixed finite
translation, so it is a homeomorphism; injectivity follows by
ample determination on $Y$. Hence its entire normalized image is
compact. If $u_j=\sigma_c+\pi^*w_j$ decreases, (13) gives
divisorial decrease of $w_j$, and their ample maximum and
determination give all-point decrease. Their limiting generic
value is that of $u_j$. Apply $G_0$ to obtain
$w=\inf_j w_j$ on $Y$. Literal addition and pullback give
$\inf_j u_j=\sigma_c+\pi^*w$ at all points; the forward
construction supplies the required one actual FS defining net.

To prove the first part of (10), the nef-intersection lemma makes
every $v\in\PSH(c_1(\pi^*A))$ genuine. Its finite FS net from
that lemma, tensored with $s_{mcE}$, gives a genuine
$\sigma_c+v$ in class $L$. Equation (9) gives
$\sigma_c+v=\sigma_c+\pi^*w$ for some $w\in\PSH(Y,A)$.
Subtract the finite exceptional values on divisorial coordinates,
and use common-ample all-point determination on $X$ to get
$v=\pi^*w$. The reverse inclusion was already constructed by
pulling back the finite ample FS net. Finally (7) for $B=A,b=c$
gives $h^0(X,mL)=h^0(Y,mA)$ on all divisible powers. The ample
Hilbert polynomial has leading coefficient $A^n/n!$, proving
the raw volume in (10).
\end{proof}

\section{What this repairs and what remains}
The theorem concerns the entire nonnef big cone in the stated
blowup family, not a selected subcone with a convenient gauge.
Its class can be nonnef in arbitrarily large dimension. The proof
retains arbitrary-cardinality nets, actual divisor powers and
all-point endpoints. Normal-section extension, not a hypothetical
normal differential of the blowup, supplies the divisibility.
This theorem is relative only to the exact ample analytic premise
and the explicit ordinary algebraic inputs above; it has not had
an independent review and adds no accepted paper or solved atom.

An arbitrary big class need not have this exceptional presentation.
That original domain still needs its own genuine closure and
compactness construction. Even in this family (9) does not by itself
construct an original-carrier positive finite-energy measure:
lifting an arbitrary measure through $\pi^{\rm an}$ requires
the correct fiber and polar-set treatment, and the genuine
intersection/operator and weighted unbounded limits must be proved.
The raw volume in (10) cannot substitute for those steps.
No intrinsic target, two-sided energy derivative, all-solution
comparison, variational attainment or uniqueness assertion is made
here. Zero-volume pseudoeffective classes remain a separate domain.
The full unchanged original P16.5 remains open in this research record.
\end{document}
